\section{Message Authentication Code (MAC)}
Fungsi hash standar menjamin integritas (bahwa pesan tidak berubah), tetapi tidak menjamin otentikasi (siapa yang mengirim pesan). Untuk itu, digunakan \textbf{MAC}.

\subsection{Konsep MAC}
MAC adalah fungsi hash yang menggunakan kunci rahasia bersama ($K$).
\[ \text{Tag} = \text{MAC}(K, m) \]
Hanya pihak yang memiliki kunci $K$ yang dapat membangkitkan atau memverifikasi tag tersebut. Hal ini memastikan bahwa pesan benar-benar berasal dari pihak yang memegang kunci.

\subsection{HMAC (Hash-based MAC)}
HMAC adalah konstruksi MAC yang menggunakan fungsi hash kriptografis (seperti SHA-256) dengan menambahkan kunci rahasia melalui proses \textit{padding} dan \textit{hashing} dua lapis untuk mencegah serangan \textit{length-extension}.

% Konstruksi HMAC: hashing dua lapis dengan kunci yang dipad.
% Dirujuk dari section-03.tex dengan % Konstruksi HMAC: hashing dua lapis dengan kunci yang dipad.
% Dirujuk dari section-03.tex dengan % Konstruksi HMAC: hashing dua lapis dengan kunci yang dipad.
% Dirujuk dari section-03.tex dengan \input{figures/skema-hmac}.
\begin{figure}[htbp]
    \centering
    \begin{tikzpicture}[
        kotak/.style={draw, rounded corners, align=center, font=\small, inner sep=5pt,
                      minimum width=5.6cm, minimum height=0.75cm},
        hash/.style={kotak, minimum width=2.2cm, fill=orange!12},
        panah/.style={-{Latex[length=2.4mm]}, thick},
        node distance=0.45cm,
    ]
        \node[kotak, fill=blue!6] (a) {$M \;\|\; (K \oplus \text{ipad})$};
        \node[hash, below=of a] (h1) {$H$};
        \node[kotak, fill=gray!8, below=of h1] (b) {$h_1 \;\|\; (K \oplus \text{opad})$};
        \node[hash, below=of b] (h2) {$H$};
        \node[kotak, fill=green!10, below=of h2] (tag) {Tag MAC};

        \draw[panah] (a) -- (h1);
        \draw[panah] (h1) -- (b);
        \draw[panah] (b) -- (h2);
        \draw[panah] (h2) -- (tag);

        \node[right=0.25cm of h1, font=\scriptsize, align=left]
            {$h_1 = H\bigl((K \oplus \text{ipad}) \,\|\, M\bigr)$};
        \node[right=0.25cm of h2, font=\scriptsize, align=left]
            {$\text{Tag} = H\bigl((K \oplus \text{opad}) \,\|\, h_1\bigr)$};
    \end{tikzpicture}
    \caption{Konstruksi HMAC: hashing dua lapis dengan kunci yang dipad memakai ipad dan opad}
    \label{fig:skema-hmac}
\end{figure}
.
\begin{figure}[htbp]
    \centering
    \begin{tikzpicture}[
        kotak/.style={draw, rounded corners, align=center, font=\small, inner sep=5pt,
                      minimum width=5.6cm, minimum height=0.75cm},
        hash/.style={kotak, minimum width=2.2cm, fill=orange!12},
        panah/.style={-{Latex[length=2.4mm]}, thick},
        node distance=0.45cm,
    ]
        \node[kotak, fill=blue!6] (a) {$M \;\|\; (K \oplus \text{ipad})$};
        \node[hash, below=of a] (h1) {$H$};
        \node[kotak, fill=gray!8, below=of h1] (b) {$h_1 \;\|\; (K \oplus \text{opad})$};
        \node[hash, below=of b] (h2) {$H$};
        \node[kotak, fill=green!10, below=of h2] (tag) {Tag MAC};

        \draw[panah] (a) -- (h1);
        \draw[panah] (h1) -- (b);
        \draw[panah] (b) -- (h2);
        \draw[panah] (h2) -- (tag);

        \node[right=0.25cm of h1, font=\scriptsize, align=left]
            {$h_1 = H\bigl((K \oplus \text{ipad}) \,\|\, M\bigr)$};
        \node[right=0.25cm of h2, font=\scriptsize, align=left]
            {$\text{Tag} = H\bigl((K \oplus \text{opad}) \,\|\, h_1\bigr)$};
    \end{tikzpicture}
    \caption{Konstruksi HMAC: hashing dua lapis dengan kunci yang dipad memakai ipad dan opad}
    \label{fig:skema-hmac}
\end{figure}
.
\begin{figure}[htbp]
    \centering
    \begin{tikzpicture}[
        kotak/.style={draw, rounded corners, align=center, font=\small, inner sep=5pt,
                      minimum width=5.6cm, minimum height=0.75cm},
        hash/.style={kotak, minimum width=2.2cm, fill=orange!12},
        panah/.style={-{Latex[length=2.4mm]}, thick},
        node distance=0.45cm,
    ]
        \node[kotak, fill=blue!6] (a) {$M \;\|\; (K \oplus \text{ipad})$};
        \node[hash, below=of a] (h1) {$H$};
        \node[kotak, fill=gray!8, below=of h1] (b) {$h_1 \;\|\; (K \oplus \text{opad})$};
        \node[hash, below=of b] (h2) {$H$};
        \node[kotak, fill=green!10, below=of h2] (tag) {Tag MAC};

        \draw[panah] (a) -- (h1);
        \draw[panah] (h1) -- (b);
        \draw[panah] (b) -- (h2);
        \draw[panah] (h2) -- (tag);

        \node[right=0.25cm of h1, font=\scriptsize, align=left]
            {$h_1 = H\bigl((K \oplus \text{ipad}) \,\|\, M\bigr)$};
        \node[right=0.25cm of h2, font=\scriptsize, align=left]
            {$\text{Tag} = H\bigl((K \oplus \text{opad}) \,\|\, h_1\bigr)$};
    \end{tikzpicture}
    \caption{Konstruksi HMAC: hashing dua lapis dengan kunci yang dipad memakai ipad dan opad}
    \label{fig:skema-hmac}
\end{figure}


\subsection{Rumus Lengkap HMAC}
HMAC didefinisikan secara eksplisit sebagai
\[ \mathrm{HMAC}(K,m) = H\big((K^{+}\oplus \mathrm{opad}) \,\|\, H((K^{+}\oplus \mathrm{ipad}) \,\|\, m)\big), \]
dengan $K^{+}$ adalah kunci $K$ yang dipad nol hingga sepanjang satu blok hash ($B$ byte: 64 untuk SHA-256, 128 untuk SHA-512), atau di-hash lebih dulu bila $|K| > B$. Konstanta $\mathrm{ipad} = \code{0x36}$ dan $\mathrm{opad} = \code{0x5C}$ diulang $B$ kali, sehingga $H$ dipanggil dua kali melalui dua jalur kunci efektif yang berbeda \citep{stallings2017,menezes1996}.

\subsection{Mengapa Bukan \texorpdfstring{$H(K\|m)$}{H(K||m)}?}
Menempelkan kunci di depan pesan, $H(K\|m)$, rentan \textbf{length extension}: karena sifat Merkle--Damgård, penyerang yang mengetahui $H(K\|m)$ dan panjang $K$ dapat menghitung $H(K\|m\|\mathit{pad}\|m')$ tanpa mengetahui $K$. HMAC menutup celah ini sebab $h_1 = H((K\oplus\mathrm{ipad})\|m)$ belum menjadi tag; tag baru terbentuk setelah hash lapis kedua dengan opad. HMAC tetap aman selama fungsi kompresi dianggap tahan \textit{pseudorandom} \citep{katz2020}.

\subsection{MAC vs Tanda Tangan Digital vs Checksum}
\begin{itemize}
    \item \textbf{Checksum/CRC}: mendeteksi galat acak, tanpa kunci. Penyerang bebas menghitung ulang nilainya, jadi integritas terhadap manipulasi tidak terjamin.
    \item \textbf{MAC/HMAC}: memakai kunci simetris bersama. Menjamin integritas dan otentikasi pesan, tetapi \emph{tidak} menjamin \textit{non-repudiation} karena kedua pihak memegang kunci yang sama.
    \item \textbf{Tanda tangan digital}: memakai kunci privat/publik. Menjamin integritas, otentikasi, dan \textit{non-repudiation}, tetapi biaya komputasinya jauh lebih tinggi daripada MAC.
\end{itemize}





